\documentclass[11pt,a4paper]{article}

% -------------------------------------------------
% Packages
% -------------------------------------------------
\usepackage[margin=2.2cm]{geometry}
\usepackage{kotex}
\usepackage{amsmath,amssymb}
\usepackage{booktabs}
\usepackage{array}
\usepackage{xcolor}
\usepackage{listings}
\usepackage[most]{tcolorbox}
\usepackage{enumitem}
\usepackage{hyperref}
\usepackage{fancyhdr}
\usepackage{titlesec}

% -------------------------------------------------
% Colors
% -------------------------------------------------
\definecolor{codebg}{RGB}{247,247,247}
\definecolor{codeframe}{RGB}{210,210,210}
\definecolor{keywordcolor}{RGB}{0,70,160}
\definecolor{commentcolor}{RGB}{0,120,80}
\definecolor{stringcolor}{RGB}{170,50,50}
\definecolor{titleblue}{RGB}{35,75,130}

% -------------------------------------------------
% Hyperlinks
% -------------------------------------------------
\hypersetup{
    colorlinks=true,
    linkcolor=titleblue,
    urlcolor=titleblue
}

% -------------------------------------------------
% Header / Footer
% -------------------------------------------------
\pagestyle{fancy}
\fancyhf{}
\lhead{C++ Programming}
\rhead{Day 4}
\cfoot{\thepage}

% -------------------------------------------------
% Section style
% -------------------------------------------------
\titleformat{\section}
  {\Large\bfseries\color{titleblue}}
  {\thesection.}{0.6em}{}

\titleformat{\subsection}
  {\large\bfseries}
  {\thesubsection.}{0.5em}{}

% -------------------------------------------------
% C++ code style
% -------------------------------------------------
\lstdefinestyle{cppstyle}{
    language=C++,
    backgroundcolor=\color{codebg},
    frame=single,
    rulecolor=\color{codeframe},
    basicstyle=\ttfamily\small,
    keywordstyle=\color{keywordcolor}\bfseries,
    commentstyle=\color{commentcolor},
    stringstyle=\color{stringcolor},
    numbers=left,
    numberstyle=\tiny\color{gray},
    numbersep=8pt,
    tabsize=4,
    showstringspaces=false,
    breaklines=true,
    columns=fullflexible
}

\lstset{style=cppstyle}

% -------------------------------------------------
% Boxes
% -------------------------------------------------
\newtcolorbox{learningbox}{
    colback=blue!4,
    colframe=titleblue,
    title=\textbf{학습 목표},
    fonttitle=\bfseries
}

\newtcolorbox{importantbox}{
    colback=yellow!8,
    colframe=orange!70!black,
    title=\textbf{핵심 포인트},
    fonttitle=\bfseries
}

\newtcolorbox{practicebox}{
    colback=red!3,
    colframe=red!55!black,
    title=\textbf{실습},
    fonttitle=\bfseries
}

% -------------------------------------------------
% Document
% -------------------------------------------------
\begin{document}

\begin{titlepage}
    \centering
    \vspace*{3cm}

    {\Huge\bfseries C++ Programming\par}
    \vspace{0.8cm}
    {\LARGE Day 4: 참조와 동적 메모리\par}

    \vspace{2cm}

    {\large Reference, Stack and Heap, Dynamic Allocation, Dynamic Array\par}

    \vfill

\end{titlepage}

\tableofcontents
\newpage

% =================================================
\section{학습 목표}
% =================================================

\begin{learningbox}
이 장을 마치면 다음 내용을 설명하고 직접 코드로 작성할 수 있다.
\begin{itemize}[leftmargin=1.5em]
    \item 참조(reference)의 개념과 별칭(alias)의 의미
    \item 값 복사와 참조 연결의 차이
    \item 참조의 초기화 규칙과 재연결 불가능성
    \item 값 전달과 참조 전달의 차이
    \item 참조를 이용한 \texttt{swap} 함수 구현
    \item 스택 메모리와 힙 메모리의 차이
    \item \texttt{new}, \texttt{delete}, \texttt{nullptr}의 사용
    \item 댕글링 포인터와 메모리 누수의 원인
    \item 동적 배열과 \texttt{new[]}, \texttt{delete[]}
    \item 배열과 포인터 연산의 관계
\end{itemize}
\end{learningbox}

% =================================================
\section{참조란 무엇인가}
% =================================================

참조(reference)는 이미 존재하는 변수에 붙이는 또 하나의 이름이다.
새로운 값을 저장하는 독립적인 변수를 만드는 것이 아니라, 원래 변수와 같은 메모리 공간을 가리키는 별칭(alias)을 만든다.

\begin{lstlisting}
#include <iostream>

using namespace std;

int main()
{
    int number = 10;
    int& ref = number;

    cout << number << endl;
    cout << ref << endl;

    return 0;
}
\end{lstlisting}

출력 결과:

\begin{verbatim}
10
10
\end{verbatim}

\texttt{ref}는 \texttt{number}의 복사본이 아니다.
두 이름은 같은 정수 객체를 나타낸다.

\subsection{참조 선언 문법}

참조는 자료형 뒤에 \texttt{\&}를 붙여 선언한다.

\begin{lstlisting}
int value = 20;
int& reference = value;
\end{lstlisting}

여기서 \texttt{int\&}는 ``정수형 변수에 대한 참조''라는 뜻이다.

\begin{importantbox}
선언문에서의 \texttt{\&}는 참조를 나타낸다.
표현식에서의 \texttt{\&variable}은 해당 변수의 주소를 구하는 주소 연산자이다.
모양은 같지만 사용되는 위치와 의미가 다르다.
\end{importantbox}

% =================================================
\section{값 복사와 참조의 차이}
% =================================================

일반 변수를 다른 변수로 초기화하면 값이 복사된다.

\begin{lstlisting}
#include <iostream>

using namespace std;

int main()
{
    int original = 10;
    int copied = original;

    copied = 50;

    cout << original << endl;
    cout << copied << endl;

    return 0;
}
\end{lstlisting}

출력 결과:

\begin{verbatim}
10
50
\end{verbatim}

\texttt{original}과 \texttt{copied}는 서로 다른 메모리 공간을 가진다.
따라서 \texttt{copied}를 수정해도 \texttt{original}은 변하지 않는다.

반면 참조는 원본과 같은 메모리 공간을 사용한다.

\begin{lstlisting}
#include <iostream>

using namespace std;

int main()
{
    int original = 10;
    int& ref = original;

    ref = 50;

    cout << original << endl;
    cout << ref << endl;

    return 0;
}
\end{lstlisting}

출력 결과:

\begin{verbatim}
50
50
\end{verbatim}

참조를 통해 값을 변경하면 원본도 함께 변경된다.
사실 두 값이 동시에 바뀐 것이 아니라 하나의 객체를 두 이름으로 접근한 것이다.

\subsection{메모리 관점에서 비교}

\begin{center}
\begin{tabular}{lll}
\toprule
구분 & 값 복사 & 참조 \\
\midrule
선언 & \texttt{int b = a;} & \texttt{int\& b = a;} \\
메모리 공간 & 서로 다름 & 동일함 \\
\texttt{b} 수정 시 \texttt{a} & 변하지 않음 & 함께 변함 \\
독립된 객체 여부 & 독립된 객체 & 원본의 별칭 \\
\bottomrule
\end{tabular}
\end{center}

주소를 출력하면 차이를 더 분명하게 확인할 수 있다.

\begin{lstlisting}
#include <iostream>

using namespace std;

int main()
{
    int original = 10;
    int copied = original;
    int& ref = original;

    cout << &original << endl;
    cout << &copied << endl;
    cout << &ref << endl;

    return 0;
}
\end{lstlisting}

\texttt{original}과 \texttt{ref}의 주소는 같고, \texttt{copied}의 주소는 다르다.

% =================================================
\section{참조의 초기화 규칙}
% =================================================

\subsection{선언과 동시에 초기화해야 한다}

참조는 반드시 선언하는 순간 연결할 변수를 지정해야 한다.

\begin{lstlisting}
int number = 10;
int& ref = number;
\end{lstlisting}

다음과 같이 참조만 선언하는 것은 허용되지 않는다.

\begin{lstlisting}
int& ref;
\end{lstlisting}

참조는 독립적인 저장 공간이 아니라 기존 객체의 별칭이므로, 처음부터 어떤 객체의 별칭인지 정해져 있어야 한다.

\subsection{참조는 다른 변수로 재연결할 수 없다}

한 번 연결된 참조는 다른 변수의 참조로 바뀌지 않는다.

\begin{lstlisting}
#include <iostream>

using namespace std;

int main()
{
    int first = 10;
    int second = 20;
    int& ref = first;

    ref = second;

    cout << first << endl;
    cout << second << endl;
    cout << ref << endl;

    return 0;
}
\end{lstlisting}

출력 결과:

\begin{verbatim}
20
20
20
\end{verbatim}

\texttt{ref = second;}는 참조를 \texttt{second}로 다시 연결하는 문장이 아니다.
\texttt{second}의 값을 \texttt{ref}가 가리키는 원본, 즉 \texttt{first}에 대입한다.

\subsection{참조와 포인터의 차이}

참조와 포인터는 모두 원본 객체에 간접적으로 접근할 수 있지만 사용 방식이 다르다.

\begin{center}
\begin{tabular}{lll}
\toprule
구분 & 참조 & 포인터 \\
\midrule
선언 & \texttt{int\& ref = value;} & \texttt{int* ptr = \&value;} \\
초기화 & 반드시 필요 & 나중에 가능 \\
다른 객체로 변경 & 불가능 & 가능 \\
널 상태 & 일반적으로 없음 & \texttt{nullptr} 가능 \\
값 접근 & \texttt{ref} & \texttt{*ptr} \\
\bottomrule
\end{tabular}
\end{center}

참조는 일반 변수처럼 간단하게 사용할 수 있고, 포인터는 주소 자체를 저장하고 변경할 수 있다는 차이가 있다.

% =================================================
\section{함수의 값 전달}
% =================================================

함수의 매개변수를 일반 변수로 선언하면 인자의 값이 복사된다.
이를 값 전달(pass by value)이라고 한다.

\begin{lstlisting}
#include <iostream>

using namespace std;

void increase(int number)
{
    number++;
}

int main()
{
    int value = 10;

    increase(value);

    cout << value << endl;

    return 0;
}
\end{lstlisting}

출력 결과:

\begin{verbatim}
10
\end{verbatim}

\texttt{increase()}의 \texttt{number}는 \texttt{value}의 복사본이다.
함수 내부에서 복사본을 수정했으므로 원본은 변하지 않는다.

\subsection{값 전달의 메모리 구조}

함수를 호출하면 매개변수는 함수 내부에서 사용하는 새로운 지역 변수로 만들어진다.

\begin{verbatim}
main의 value:      10
increase의 number: 10
\end{verbatim}

두 변수의 초기값은 같지만 서로 다른 객체이다.

% =================================================
\section{참조 전달}
% =================================================

매개변수를 참조로 선언하면 인자의 복사본을 만들지 않고 원본에 직접 접근한다.
이를 참조 전달(pass by reference)이라고 한다.

\begin{lstlisting}
#include <iostream>

using namespace std;

void increase(int& number)
{
    number++;
}

int main()
{
    int value = 10;

    increase(value);

    cout << value << endl;

    return 0;
}
\end{lstlisting}

출력 결과:

\begin{verbatim}
11
\end{verbatim}

함수의 \texttt{number}는 \texttt{main()}의 \texttt{value}에 대한 참조이다.
따라서 함수 안의 변경이 원본에 반영된다.

\subsection{값 전달과 참조 전달 비교}

\begin{center}
\begin{tabular}{lll}
\toprule
구분 & 값 전달 & 참조 전달 \\
\midrule
매개변수 & \texttt{int number} & \texttt{int\& number} \\
인자 복사 & 발생 & 발생하지 않음 \\
원본 수정 & 불가능 & 가능 \\
큰 객체 전달 비용 & 상대적으로 큼 & 상대적으로 작음 \\
\bottomrule
\end{tabular}
\end{center}

\begin{importantbox}
참조 전달은 원본을 수정해야 할 때 유용하다.
다만 함수 호출만 보고 원본이 바뀌는지 알기 어려울 수 있으므로, 함수 이름과 역할을 명확하게 정하는 것이 중요하다.
\end{importantbox}

% =================================================
\section{참조를 이용한 \texttt{swap}}
% =================================================

두 변수의 값을 바꾸는 함수를 값 전달로 작성하면 원하는 결과를 얻을 수 없다.

\begin{lstlisting}
void wrongSwap(int a, int b)
{
    int temp = a;
    a = b;
    b = temp;
}
\end{lstlisting}

위 함수에서 바뀌는 것은 매개변수 \texttt{a}, \texttt{b}뿐이다.
호출한 쪽의 원본 변수는 변하지 않는다.

참조 전달을 사용하면 원본 변수의 값을 직접 바꿀 수 있다.

\begin{lstlisting}
#include <iostream>

using namespace std;

void swapValues(int& a, int& b)
{
    int temp = a;
    a = b;
    b = temp;
}

int main()
{
    int x = 10;
    int y = 20;

    swapValues(x, y);

    cout << x << endl;
    cout << y << endl;

    return 0;
}
\end{lstlisting}

출력 결과:

\begin{verbatim}
20
10
\end{verbatim}

호출 시점에는 일반 변수처럼 \texttt{swapValues(x, y)}라고 작성하지만, 함수 정의에서 매개변수가 참조로 선언되어 있으므로 \texttt{a}와 \texttt{b}는 각각 \texttt{x}와 \texttt{y}의 별칭이 된다.

% =================================================
\section{스택 메모리와 힙 메모리}
% =================================================

프로그램이 실행되면 필요한 데이터가 메모리에 저장된다.
C++에서 지역 변수와 동적 할당을 이해하려면 스택(stack)과 힙(heap)의 차이를 알아야 한다.

\subsection{스택 메모리}

함수 내부에서 선언한 지역 변수는 일반적으로 스택 메모리에 저장된다.

\begin{lstlisting}
void function()
{
    int number = 10;
    double value = 3.14;
}
\end{lstlisting}

함수가 호출되면 지역 변수가 만들어지고, 함수가 끝나면 자동으로 제거된다.
프로그래머가 직접 메모리를 해제할 필요가 없다.

\subsection{힙 메모리}

힙은 프로그램 실행 중 필요한 크기만큼 직접 메모리를 요청할 때 사용하는 영역이다.
C++에서는 \texttt{new} 연산자로 힙 메모리를 할당하고, \texttt{delete} 연산자로 해제한다.

\begin{center}
\begin{tabular}{lll}
\toprule
구분 & 스택 & 힙 \\
\midrule
생성 방식 & 지역 변수 선언 & \texttt{new} 사용 \\
해제 방식 & 자동 & \texttt{delete} 필요 \\
수명 & 블록 또는 함수가 끝날 때까지 & 직접 해제할 때까지 \\
관리 주체 & 컴파일러와 실행 환경 & 프로그래머 \\
주요 위험 & 제한된 크기 & 누수, 댕글링 포인터 \\
\bottomrule
\end{tabular}
\end{center}

\begin{importantbox}
스택과 힙의 핵심 차이는 ``자동 관리''와 ``직접 관리''이다.
힙 메모리는 필요한 시점과 크기를 유연하게 정할 수 있지만, 반드시 올바르게 해제해야 한다.
\end{importantbox}

% =================================================
\section{동적 메모리 할당: \texttt{new}}
% =================================================

\texttt{new}는 힙에 메모리를 할당하고, 할당된 메모리의 주소를 반환한다.
따라서 반환된 주소를 저장할 포인터가 필요하다.

\begin{lstlisting}
int* ptr = new int;
\end{lstlisting}

위 코드는 힙에 정수 하나를 저장할 공간을 만들고, 그 주소를 \texttt{ptr}에 저장한다.

\subsection{동적으로 할당한 값 사용}

포인터가 가리키는 값에 접근하려면 역참조 연산자 \texttt{*}를 사용한다.

\begin{lstlisting}
#include <iostream>

using namespace std;

int main()
{
    int* ptr = new int;

    *ptr = 42;

    cout << *ptr << endl;

    delete ptr;

    return 0;
}
\end{lstlisting}

\texttt{ptr} 자체는 주소를 저장하는 포인터 변수이고, \texttt{*ptr}은 힙에 실제로 저장된 정수이다.

\subsection{할당과 동시에 초기화}

\begin{lstlisting}
int* ptr = new int(42);
\end{lstlisting}

중괄호 초기화도 사용할 수 있다.

\begin{lstlisting}
int* ptr = new int{42};
\end{lstlisting}

두 경우 모두 정수 공간을 만들고 초기값 42를 저장한다.

% =================================================
\section{동적 메모리 해제: \texttt{delete}}
% =================================================

\texttt{new}로 할당한 메모리는 사용이 끝난 뒤 \texttt{delete}로 해제해야 한다.

\begin{lstlisting}
int* ptr = new int(42);

cout << *ptr << endl;

delete ptr;
\end{lstlisting}

\texttt{delete ptr;}는 \texttt{ptr} 변수를 제거하는 것이 아니다.
\texttt{ptr}이 가리키는 힙 메모리를 해제한다.
포인터 변수 자체는 해당 블록이 끝날 때까지 남아 있다.

\subsection{해제 후 포인터 상태}

메모리를 해제한 뒤에도 포인터에는 이전 주소가 남아 있다.
그러나 그 주소의 메모리는 더 이상 사용할 수 없다.

\begin{lstlisting}
int* ptr = new int(42);

delete ptr;

cout << *ptr << endl;
\end{lstlisting}

위 코드의 마지막 줄은 정의되지 않은 동작(undefined behavior)을 일으킨다.
프로그램이 우연히 값을 출력할 수도 있고, 오류가 발생할 수도 있다.

% =================================================
\section{\texttt{nullptr}}
% =================================================

\texttt{nullptr}는 포인터가 유효한 객체를 가리키고 있지 않음을 나타내는 값이다.

\begin{lstlisting}
int* ptr = nullptr;
\end{lstlisting}

포인터를 아직 어떤 주소와도 연결하지 않을 때 초기값으로 사용할 수 있다.

\begin{lstlisting}
#include <iostream>

using namespace std;

int main()
{
    int* ptr = nullptr;

    if (ptr == nullptr)
    {
        cout << "No object" << endl;
    }

    return 0;
}
\end{lstlisting}

메모리를 해제한 뒤 포인터를 \texttt{nullptr}로 바꾸면 잘못된 주소를 다시 사용하는 실수를 줄일 수 있다.

\begin{lstlisting}
int* ptr = new int(42);

delete ptr;
ptr = nullptr;
\end{lstlisting}

\begin{importantbox}
\texttt{delete nullptr;}는 안전하다.
포인터가 \texttt{nullptr}인지 검사한 뒤에만 삭제해야 하는 것은 아니지만, 포인터가 현재 유효한 메모리를 가리키는지는 항상 구분해야 한다.
\end{importantbox}

% =================================================
\section{댕글링 포인터}
% =================================================

댕글링 포인터(dangling pointer)는 이미 수명이 끝났거나 해제된 메모리의 주소를 계속 저장하고 있는 포인터이다.

\begin{lstlisting}
int* ptr = new int(42);

delete ptr;
\end{lstlisting}

\texttt{delete} 이후 \texttt{ptr}은 더 이상 유효한 정수를 가리키지 않지만 이전 주소는 그대로 가지고 있다.
이 상태의 \texttt{ptr}이 댕글링 포인터이다.

\subsection{댕글링 포인터 사용의 위험}

\begin{lstlisting}
int* ptr = new int(42);

delete ptr;

*ptr = 100;
\end{lstlisting}

이미 해제된 메모리에 값을 쓰려고 하므로 정의되지 않은 동작이 발생한다.
다른 데이터가 손상되거나 프로그램이 비정상 종료될 수 있다.

\subsection{예방 방법}

가장 기본적인 방법은 해제 직후 포인터를 \texttt{nullptr}로 만드는 것이다.

\begin{lstlisting}
delete ptr;
ptr = nullptr;
\end{lstlisting}

또한 여러 포인터가 같은 동적 메모리를 가리킬 때는 한 포인터가 메모리를 해제한 뒤 다른 포인터도 모두 무효가 된다는 점을 주의해야 한다.

\begin{lstlisting}
int* first = new int(42);
int* second = first;

delete first;
first = nullptr;
\end{lstlisting}

이때 \texttt{second}는 여전히 이전 주소를 저장하고 있으므로 댕글링 포인터이다.

% =================================================
\section{메모리 누수}
% =================================================

메모리 누수(memory leak)는 동적으로 할당한 메모리를 더 이상 사용할 방법이 없는데도 해제하지 않은 상태를 말한다.

\begin{lstlisting}
void createNumber()
{
    int* ptr = new int(42);
}
\end{lstlisting}

함수가 끝나면 지역 포인터 \texttt{ptr}은 사라진다.
하지만 \texttt{new}로 만든 힙 메모리는 자동으로 해제되지 않는다.
그 주소를 잃어버렸으므로 이후에는 \texttt{delete}할 수도 없다.

\subsection{주소를 덮어쓰는 경우}

\begin{lstlisting}
int* ptr = new int(10);

ptr = new int(20);
\end{lstlisting}

두 번째 대입이 실행되면 첫 번째로 할당한 정수의 주소를 잃어버린다.
첫 번째 메모리는 프로그램이 종료될 때까지 남아 메모리 누수가 된다.

올바른 순서는 기존 메모리를 먼저 해제한 뒤 새로운 메모리를 할당하는 것이다.

\begin{lstlisting}
int* ptr = new int(10);

delete ptr;
ptr = new int(20);

delete ptr;
ptr = nullptr;
\end{lstlisting}

\subsection{반복문 안의 메모리 누수}

작은 누수도 반복되면 큰 문제가 된다.

\begin{lstlisting}
for (int i = 0; i < 1000000; i++)
{
    int* ptr = new int(i);
}
\end{lstlisting}

각 반복에서 포인터 변수는 사라지지만 할당된 정수는 해제되지 않는다.
반복문이 진행될수록 사용 중인 메모리가 계속 증가한다.

% =================================================
\section{동적 배열}
% =================================================

일반 배열의 크기는 선언 시점에 정해져야 한다.
실행 중 입력받은 크기로 배열을 만들고 싶을 때 동적 배열을 사용할 수 있다.

\begin{lstlisting}
int size;
cin >> size;

int* numbers = new int[size];
\end{lstlisting}

\texttt{new int[size]}는 정수 \texttt{size}개가 연속해서 저장될 공간을 힙에 만들고 첫 번째 원소의 주소를 반환한다.

\subsection{동적 배열 원소 접근}

동적 배열도 일반 배열과 같은 대괄호 문법으로 접근할 수 있다.

\begin{lstlisting}
#include <iostream>

using namespace std;

int main()
{
    int size;

    cout << "Size: ";
    cin >> size;

    int* numbers = new int[size];

    for (int i = 0; i < size; i++)
    {
        numbers[i] = (i + 1) * 10;
    }

    for (int i = 0; i < size; i++)
    {
        cout << numbers[i] << ' ';
    }

    cout << endl;

    delete[] numbers;
    numbers = nullptr;

    return 0;
}
\end{lstlisting}

\subsection{\texttt{delete[]}로 해제}

\texttt{new[]}로 만든 배열은 반드시 \texttt{delete[]}로 해제해야 한다.

\begin{lstlisting}
int* numbers = new int[5];

delete[] numbers;
numbers = nullptr;
\end{lstlisting}

단일 객체와 배열의 할당 및 해제 방식은 짝을 맞춰야 한다.

\begin{center}
\begin{tabular}{lll}
\toprule
대상 & 할당 & 해제 \\
\midrule
단일 객체 & \texttt{new int} & \texttt{delete ptr} \\
배열 & \texttt{new int[n]} & \texttt{delete[] ptr} \\
\bottomrule
\end{tabular}
\end{center}

\begin{importantbox}
\texttt{new[]}에 \texttt{delete}를 사용하거나, 단일 객체에 \texttt{delete[]}를 사용하는 것은 잘못된 코드이다.
할당 방식과 해제 방식은 반드시 정확히 대응해야 한다.
\end{importantbox}

% =================================================
\section{동적 배열 초기화와 입력}
% =================================================

동적 배열도 생성과 동시에 값 초기화를 할 수 있다.

\begin{lstlisting}
int* numbers = new int[5]{};
\end{lstlisting}

빈 중괄호를 사용하면 모든 정수 원소가 0으로 초기화된다.

일부 값을 직접 지정할 수도 있다.

\begin{lstlisting}
int* numbers = new int[5]{10, 20, 30, 40, 50};
\end{lstlisting}

사용자로부터 원소를 입력받는 예제는 다음과 같다.

\begin{lstlisting}
#include <iostream>

using namespace std;

int main()
{
    int size;
    int sum = 0;

    cout << "Size: ";
    cin >> size;

    int* numbers = new int[size];

    for (int i = 0; i < size; i++)
    {
        cout << "numbers[" << i << "]: ";
        cin >> numbers[i];
        sum += numbers[i];
    }

    cout << "Sum: " << sum << endl;

    delete[] numbers;
    numbers = nullptr;

    return 0;
}
\end{lstlisting}

동적 배열은 사용이 끝난 지점에서 즉시 해제하는 습관이 중요하다.

% =================================================
\section{배열과 포인터의 관계}
% =================================================

배열 이름은 대부분의 표현식에서 첫 번째 원소의 주소처럼 사용된다.

\begin{lstlisting}
int numbers[5] = {10, 20, 30, 40, 50};

cout << numbers << endl;
cout << &numbers[0] << endl;
\end{lstlisting}

두 출력은 같은 주소를 나타낸다.

\subsection{인덱스와 역참조}

배열의 인덱스 표현은 포인터 연산으로 바꾸어 쓸 수 있다.

\begin{lstlisting}
numbers[i]
*(numbers + i)
\end{lstlisting}

두 표현은 같은 원소를 나타낸다.

\begin{lstlisting}
#include <iostream>

using namespace std;

int main()
{
    int numbers[5] = {10, 20, 30, 40, 50};

    cout << numbers[2] << endl;
    cout << *(numbers + 2) << endl;

    return 0;
}
\end{lstlisting}

출력 결과:

\begin{verbatim}
30
30
\end{verbatim}

\subsection{주소가 증가하는 단위}

정수 포인터에 1을 더하면 주소가 1바이트 증가하는 것이 아니라 다음 정수의 주소로 이동한다.

\begin{lstlisting}
int numbers[3] = {10, 20, 30};
int* ptr = numbers;

cout << *ptr << endl;
ptr++;
cout << *ptr << endl;
\end{lstlisting}

첫 번째 출력은 10이고, \texttt{ptr++} 이후 두 번째 출력은 20이다.
포인터 연산은 가리키는 자료형의 크기를 기준으로 이동한다.

% =================================================
\section{포인터 연산으로 배열 순회}
% =================================================

배열을 인덱스로 순회하는 대신 포인터를 이동시키며 순회할 수 있다.

\begin{lstlisting}
#include <iostream>

using namespace std;

int main()
{
    int numbers[5] = {10, 20, 30, 40, 50};
    int* ptr = numbers;

    for (int i = 0; i < 5; i++)
    {
        cout << *ptr << ' ';
        ptr++;
    }

    cout << endl;

    return 0;
}
\end{lstlisting}

동적 배열에도 같은 원리가 적용된다.

\begin{lstlisting}
int* numbers = new int[5]{10, 20, 30, 40, 50};

for (int i = 0; i < 5; i++)
{
    cout << *(numbers + i) << ' ';
}

delete[] numbers;
numbers = nullptr;
\end{lstlisting}

\subsection{배열 범위를 벗어나면 안 된다}

포인터 연산으로 배열 범위를 넘어선 위치를 역참조하면 정의되지 않은 동작이 발생한다.

\begin{lstlisting}
int numbers[3] = {10, 20, 30};

cout << *(numbers + 3) << endl;
\end{lstlisting}

유효한 인덱스는 0, 1, 2뿐이다.
\texttt{numbers + 3}은 배열의 마지막 원소 다음 위치를 가리키며, 그 위치를 역참조해서는 안 된다.

% =================================================
\section{동적 메모리를 함수에서 사용하기}
% =================================================

함수 안에서 동적 배열을 만들고 주소를 반환할 수 있다.

\begin{lstlisting}
int* createArray(int size)
{
    int* numbers = new int[size]{};
    return numbers;
}
\end{lstlisting}

함수 내부의 지역 포인터 \texttt{numbers}는 함수가 끝나면 사라지지만, 힙에 생성한 배열은 자동으로 사라지지 않는다.
주소가 호출한 쪽으로 반환되므로 계속 사용할 수 있다.

\begin{lstlisting}
#include <iostream>

using namespace std;

int* createArray(int size)
{
    int* numbers = new int[size]{};
    return numbers;
}

int main()
{
    int size = 5;
    int* values = createArray(size);

    for (int i = 0; i < size; i++)
    {
        values[i] = i + 1;
    }

    for (int i = 0; i < size; i++)
    {
        cout << values[i] << ' ';
    }

    cout << endl;

    delete[] values;
    values = nullptr;

    return 0;
}
\end{lstlisting}

메모리를 생성한 함수와 해제하는 함수가 다를 수 있으므로, 어떤 코드가 해제 책임을 가지는지 분명하게 정해야 한다.

% =================================================
\section{동적 메모리 사용 원칙}
% =================================================

동적 메모리를 안전하게 사용하려면 다음 순서를 습관화해야 한다.

\begin{enumerate}
    \item 포인터를 선언하고 필요하면 \texttt{nullptr}로 초기화한다.
    \item \texttt{new} 또는 \texttt{new[]}로 메모리를 할당한다.
    \item 포인터를 통해 데이터를 사용한다.
    \item 더 이상 필요하지 않을 때 \texttt{delete} 또는 \texttt{delete[]}로 해제한다.
    \item 해제한 포인터를 \texttt{nullptr}로 바꾼다.
\end{enumerate}

\begin{lstlisting}
int* ptr = nullptr;

ptr = new int(42);

cout << *ptr << endl;

delete ptr;
ptr = nullptr;
\end{lstlisting}

동적 배열도 같은 흐름을 따른다.

\begin{lstlisting}
int* numbers = nullptr;

numbers = new int[5]{};

for (int i = 0; i < 5; i++)
{
    numbers[i] = i * 10;
}

delete[] numbers;
numbers = nullptr;
\end{lstlisting}

% =================================================
\section{실습 목록}
% =================================================

\begin{practicebox}
Day 4 실습 파일은 다음과 같이 관리한다.
\begin{itemize}[leftmargin=1.5em]
    \item \texttt{practice/p01\_reference.cpp}
    \item \texttt{practice/p02\_dynamic.cpp}
    \item \texttt{practice/p03\_combined.cpp}
\end{itemize}
\end{practicebox}

% =================================================
\section{핵심 요약}
% =================================================

\begin{enumerate}
    \item 참조는 기존 변수에 붙이는 별칭이며 원본과 같은 객체를 나타낸다.
    \item 참조는 선언과 동시에 초기화해야 하고, 이후 다른 변수로 재연결할 수 없다.
    \item 값 전달은 인자를 복사하고, 참조 전달은 함수가 원본에 직접 접근하게 한다.
    \item 두 원본 변수의 값을 바꾸려면 \texttt{swap} 함수의 매개변수를 참조로 선언해야 한다.
    \item 지역 변수는 일반적으로 스택에 저장되고 수명이 끝나면 자동으로 제거된다.
    \item \texttt{new}로 만든 힙 메모리는 \texttt{delete}로 직접 해제해야 한다.
    \item 해제된 메모리를 가리키는 포인터는 댕글링 포인터가 되며 사용해서는 안 된다.
    \item 동적 메모리의 주소를 잃어버리고 해제하지 않으면 메모리 누수가 발생한다.
    \item \texttt{new[]}로 만든 배열은 반드시 \texttt{delete[]}로 해제해야 한다.
    \item 배열의 \texttt{arr[i]}는 포인터 표현 \texttt{*(arr + i)}와 같은 의미이다.
\end{enumerate}

\end{document}
